\section{Классификация методов восстановления глубины}

На основе проведенного анализа предметной области были определены несколько ключевых критериев сопоставления рассмотренных методов.

\begin{enumerate}[label=\arabic*.]
    \item \textbf{Способ получения данных} --- активный или пассивный подход к сбору информации.

    Предложенный критерий учитывает подход к сбору исходных данных.
    Активные методы, такие как лазерное сканирование, требуют использования специализированного оборудования, что
    предполагает проведение исследований в лабораторных условиях.
    Пассивные методы, напротив, работают с доступными изображениями.
    Применение данного критерия позволяет своевременно
    оценить необходимые условия получения исходных изображений.

    \item \textbf{Дополнительная информация} --- данные о положении камер, временных метках и других параметрах.

    Необходимость в дополнительной информации ограничивает область применения метода и
    может предполагать наличие дополнительных компетенций для их использования, например, при настройке параметров камеры.

    \item \textbf{Результат} --- итоговый объект реконструкции (карта глубин, информация о структуре объекта).

    Базовым результатом является карта глубин, однако некоторые методы предоставляют дополнительную информацию:
    форма объекта, нормали к поверхностям.
    Эти дополнительные данные могут быть критически важны для полноценной трёхмерной реконструкции сцены.

    \item \textbf{Масштаб сцены} --- локальная или глобальная реконструкция.

    Масштаб сцены определяет применимость метода к задачам разного уровня: от реконструкции небольших объектов до глобальных ландшафтов.

    \item \textbf{Уровень автоматизации} --- степень автоматизированности процесса реконструкции: от ручных до автоматических методов.

    Уровень автоматизации влияет на затраты времени и ресурсов.
    Автоматизированные методы предпочтительны для задач, требующих быстрой обработки больших объемов данных,
    тогда как ручные обеспечивают больше контроля в специфических сценариях.
\end{enumerate}

Данные критерии позволяют выбрать наиболее подходящий метод восстановления глубины в зависимости от поставленной задачи,
доступных возможностей и особенностей предметной области.

В таблице~\ref{tab:methods_comparison} приведена классификация методов восстановления глубины.
\begin{table}[H]
    \centering
    \begin{tabular}{|p{2.4cm}|p{2.2cm}|p{2.6cm}|p{2.2cm}|p{2.3cm}|p{3.1cm}|}
        \hline
        Метод & Способ \newline получения данных & Доп. \newline информация & Результат & Масштаб сцены & Уровень \newline автоматизации \\
\noalign{\hrule height 2pt}
        Метод структурного освещения & Активный & Проекция шаблона & Карта глубин, форма объекта  & Локальные & Ручные \\
        \hline
        DFD & Активный или \newline пассивный & Параметры настройки камеры & Карта глубин & Локальные & Автономные \\
        \hline
        Shape from Shading (SFS) & Пассивный & Параметры освещения & Карта глубин, нормали поверхности & Локальные & Полу-автономные \\
        \hline
        Comfortable depth generation & Пассивный & - & Карта глубин & Локальные или \newline глобальные & Автономные \\
        \hline
        MiDaS & Пассивный & - & Карта глубин & Локальные или \newline глобальные & Автономные \\
        \hline
    \end{tabular}
    \caption{Классификация методов восстановления глубины из монокулярных изображений}
    \label{tab:methods_comparison}
\end{table}

\section{Выбор метода восстановления глубины архивных фотографий}

При выборе метода для восстановления глубины из старинных фотографий необходимо учитывать их специфические особенности и ограничения.

\begin{enumerate}
    \item \textbf{Ограниченность исходных данных}. Для архивных фотографий доступны только одиночные изображения, отсутствует возможность получения стереопар или последовательностей с разных ракурсов.

    \item \textbf{Отсутствие метаданных}. Архивные фотографии не содержат информации о параметрах съемки, калибровке камеры или условиях освещения.

    \item \textbf{Качество изображений}. Старинные фотографии часто имеют низкое качество, шумы, артефакты, выцветание и повреждения, что затрудняет автоматическую обработку.
\end{enumerate}

Исходя из представленной классификации, рассмотрим применимость каждого метода.

\textbf{Метод структурированного освещения} и \textbf{DFD} требуют активного вмешательства при съемке или специальных параметров камеры, что делает их неприменимыми для архивных фотографий.

\textbf{Shape from Shading} работает с моноизображениями и может восстанавливать форму объектов на основе анализа теней и освещения.
Однако метод имеет существенные ограничения: необходимость в параметрах освещения, которые невозможно достоверно определить для архивных фотографий, и применимость преимущественно к локальным сценам с четко выраженными текстурами и тенями.
Тем не менее, возможна разработка адаптированного алгоритма, в котором параметры освещения оцениваются на основе самого изображения, а область применения включает более сложные сцены за счет сегментации изображения.

\textbf{Comfortable depth generation} представляется подходящим методом, поскольку он работает с моноизображениями и не требует дополнительной информации.
Однако, для архивных фотографий с их специфическими искажениями нейросетевые подходы предположительно демонстрируют более высокую устойчивость и точность.

\textbf{Нейросетевые методы}, в частности MiDaS, представляются наиболее подходящим решением, поскольку они:
\begin{itemize}
    \item работают с одиночными изображениями,
    \item не требуют дополнительной информации о сцене или камере,
    \item обеспечивают полностью автоматическую обработку,
    \item демонстрируют высокую точность на разнообразных сценах благодаря обучению на больших наборах данных,
    \item эффективно работают как с локальными, так и с глобальными сценами, хотя качество детализации при работе с глобальными сценами снижается.
\end{itemize}

Таким образом, для повышения качества результатов разумно использовать гибридный подход:
комбинировать \textbf{нейросетевой метод} (MiDaS) для получения базовой карты глубин с элементами \textbf{Shape from Shading} для улучшения детализации в областях с выраженной текстурой и тенями.

При этом необходимо учитывать особенности архивных фотографий, реализовав предварительную обработку для устранения артефактов и повреждений, снижения шума, повышения контраста, восстановления насыщенности и других параметров изображения.

\section{Цель и постановка задачи}

Целью данной работы является разработка метода восстановления трехмерной глубины из архивных фотографий путем создания гибридного алгоритма, объединяющего преимущества нейросетевого метода и Shape from Shading, с учетом специфических особенностей исторических изображений.

Формально задачу восстановления карты глубины из архивной фотографии можно представить как функциональную модель:

\begin{equation}
D = F(I, P),
\end{equation}

где:
\begin{itemize}
    \item $D : \mathbb{R}^2 \to \mathbb{R}$ --- искомая карта глубин, сопоставляющая каждой точке $(x, y)$ изображения значение глубины;
    \item $I : \mathbb{R}^2 \to \mathbb{R}^3$ --- исходное изображение, представленное как функция, сопоставляющая каждой точке $(x, y)$ цвет в формате RGB;
    \item $P$ --- набор параметров алгоритма;
    \item $F$ --- оператор преобразования, реализующий гибридный метод.
\end{itemize}

Оператор $F$ представляет собой композицию следующих функций:

\begin{equation}
F = F_{\text{post}} \circ F_{\text{hybrid}} \circ F_{\text{pre}},
\end{equation}

где:
\begin{itemize}
    \item $F_{\text{pre}} : \mathbb{R}^2 \times \mathbb{R}^3 \to \mathbb{R}^2 \times \mathbb{R}^3$ --- функция предобработки, выполняющая шумоподавление, коррекцию яркости и контраста;
    \item $F_{\text{hybrid}} : \mathbb{R}^2 \times \mathbb{R}^3 \to \mathbb{R}^2 \times \mathbb{R}$ --- гибридная функция, объединяющая методы MiDaS и SFS;
    \item $F_{\text{post}} : \mathbb{R}^2 \times \mathbb{R} \to \mathbb{R}^2 \times \mathbb{R}$ --- функция постобработки, выполняющая сглаживание и нормализацию карты глубин.
\end{itemize}

Гибридная функция $F_{\text{hybrid}}$ определяется как:

\begin{equation} F_{\text{hybrid}}(I) = \alpha_{map} \cdot F_{\text{MiDaS}}(I) + (1 - \alpha_{map}) \cdot F_{\text{SFS}}(I), \end{equation}

где: \begin{itemize} \item $F_{\text{MiDaS}}$ --- базовая карта глубин, полученная методом MiDaS;
\item $F_{\text{SFS}}$ --- карта деталей, полученная методом Shape from Shading;
\item $w_{detail}(x,y)$ --- карта весов деталей, определяющая их значимость;
\item $\alpha_{map}(x,y)$ --- адаптивный весовой коэффициент, уменьшающийся в областях с высокой детализацией.
\end{itemize}

\section*{Вывод}

В аналитической части работы проведен обзор существующих способов восстановления глубины по единственной фотографии, а именно
методы структурного освещения, Depth from defocus, Shape from Shading, Comfortable Depth Generation и нейросетевые методы, такие как MiDaS.

Рассмотренные подходы классифицированы по ключевым критериям:
способу получения данных, необходимости в дополнительной информации, результатам реконструкции, масштабу сцены и уровню автоматизации.

Проведенный анализ позволяет сделать вывод, что оптимальным решением для восстановления глубины из архивных фотографий является гибридный метод,
объединяющий преимущества нейросетевого подхода (MiDaS) и классического метода Shape from Shading.
Такой подход обеспечивает как общую точность восстановления геометрии сцены,
так и детализацию отдельных элементов, сохраняя при этом устойчивость к дефектам, характерным для исторических изображений.

Представленное формализованное описание задач в виде функциональной модели, включающей блоки предобработки,
гибридного восстановления и постобработки, формируют теоретическую основу для разработки практического решения в следующих разделах работы.

